% Title page
\thispagestyle{empty}
\begin{titlepage}
\setcounter{page}{1}
\begin{center}

\begin{figure}[H]
\centering
\includegraphics[width = 0.8\textwidth]{Images/usyd.jpg}
\end{figure}

\doublespacing
\textsc{ \LARGE Automatic Generation of Motion Primitives for Dynamic Walking Robots}

\vspace{2cm}

\singlespacing
A thesis submitted in partial fulfilment of the requirements for the degree of \\ Bachelor of Engineering (Honours) and Bachelor of Science

\vspace{2cm}

\large
Stephen Wardrop

\vspace{0.5cm}
\normalsize
Supervisor: Dr Ian Manchester

\vfill

October, 2014

\end{center}
\end{titlepage}
\pagebreak
~
\thispagestyle{empty}
\pagenumbering{roman}
\setcounter{page}{2}
\pagebreak
\begin{abstract}
\thispagestyle{plain}
\onehalfspacing
\pagenumbering{roman}
\setcounter{page}{3}
There is significant technical and aesthetic interest in developing bipedal robots which locomote stably, efficiently and reliably over uneven terrain. The most efficient walking is achieved when the dynamics of the robot are utilised, rather than by overpowering the natural motion of the robot through the use of control input. A class of robotic walkers for which the dynamics must evolve in this manner are so-called underactuated dynamic walkers; robots for which some degrees of freedom are not subject to direct control.

Motion planning and control for this class of robots continues to be challenging due to its computational intractability, highly complex nonlinear dynamics and intrinsic instability. Some progress has been made toward achieving stable walking gaits for underactuated walkers by restricting the problem to particular domains or motions to overcome the intractability of the general problem. One such method is to prepare a library of motion primitives, which limits the on-line computation to the choice of a sequence of primitives from within the library.

Virtual constraints are a particularly useful class of primitives as they allow significant simplification of the computation required for real-time path planning of dynamic walking over uneven terrain. The contribution of this work is a general method for the production of a library of motion primitives defined as virtual constraints that sufficiently spans the feasible motions of the robot along with improvements to the method for intelligently selecting from them.
\end{abstract}

\pagebreak

\renewcommand{\abstractname}{Statement of student contribution}
\begin{abstract}
\thispagestyle{plain}
\pagenumbering{roman}
\setcounter{page}{4}
\begin{itemize}
	\item I completed the literature review largely independently, with some papers being recommended to me by my supervisor
	\item I designed the graphical interfaces for the investigation of virtual constraints
	\item I designed the method by which the motion primitives could be automatically generated
	\item I designed and implemented the simulator
	\item I incrementally extended the algorithm designed by my supervisor and Jack Umenberger \cite{manchester13planning}
	\item I derived (though not necessarily pioneered) the equations presented, except where cited
\end{itemize}
~\\~\\

The above represents an accurate summary of the student's contribution.  \\~\\~\\

Signed \\~\\~\\

\hspace{10mm}_______________ (student) \hspace{15mm} _______________ (supervisor)
\end{abstract}

\pagebreak

\renewcommand{\abstractname}{Acknowledgements}
\begin{abstract}
\thispagestyle{plain}
\pagenumbering{roman}
\onehalfspacing
\setcounter{page}{5}
I would like to thank my supervisor, Dr Ian Manchester, for his time and forbearance in helping me grasp the matter of this field. Along with him, I thank Jack Umenberger for his helpful comments and insights to get me up to speed with his previous work in this matter.

I thank <everyone who does this> for sacrificing your time and effort in the laborious proofreading of this thesis and for all of the advice and support given.

I thank Trevor Anderson for being a willing and effective co-worker on so many assignments and projects over the course of our degrees, particularly for the last couple of years. I cannot calculate the effect that your friendship and hard work have had on myself and my grades.

I extend my sincere and heartfelt thanks to my parents for their belief in me and for investing in my education from long before I knew that it was important.

I also thank my wife Laura for enduring the countless nights and weekends spent with a preoccupied or absent husband even while taking care of our young daughter. Thank you for boosting my motivation when it was lacking and for doing everything you could to give me the space I needed to finish this work. You're amazing and your encouragement and patience have been vital to me.
\end{abstract}


\pagenumbering{roman}
\setcounter{page}{6}
\tableofcontents
\clearpage
\listoffigures
\clearpage
\listoftables
\listofalgorithms
\renewcommand{\nomname}{Abbreviations and Nomenclature}
\printnomenclature
\clearpage